\section{Introduction}
Audio temporal grounding maps a natural-language query to the intervals that answer it, beyond fixed event classes~\cite{xu2021grounding,munakata2025amr}. SpotSound equips an audio-language model with interleaved timestamps and training for absent events~\cite{sun2026spotsound}. We ask a complementary question: when several real windows contain the same event categories, does the model select the one specified by the query's relation?

Consider a recording containing ``dog then bell'' and, later, ``bell then dog.'' Both descriptions are true. Reversing the query should change the answer window while the audio and event vocabulary stay fixed. A model can fail without hallucinating a sound: it returns a real occurrence answering the wrong relation. We call this a \emph{query--window binding error}. Correct event recognition alone does not determine which occurrence to return.

T-CLAP uses concatenated sounds, reversed descriptions, and temporal contrastive training~\cite{yuan2024tclap}; CompA studies event order and attribute binding through paired audio--caption matching~\cite{ghosh2024compa}. CoSTALA develops hierarchical spatio-temporal alignment~\cite{ren2026costala}, and SHINE uses compositional hard negative queries in video grounding~\cite{cheng2024shine}. Our specific setting is \emph{two inverse descriptions simultaneously true in one recording, but identifying different local timestamp answers}. The wrong window is an answer negative, not evidence that the query is globally false: supervision must distinguish query-specific answers among real windows.

We introduce RelTwin: the same sound excerpts form two opposite-order windows, and each inverse query is scored against both timestamp answers. Ordinary candidate cross-entropy augments correct-answer sequence supervision. Our contribution is this locally confusable answer construction and its implementation in generative grounding, rather than a new contrastive principle. Inference remains ordinary timestamp generation, without a candidate scorer or target-window access.

We test whether this supervision improves over same-data SFT, whether more elaborate objectives are necessary, and whether the advantage survives a timing-cue stress test. Paired diagnostics check that answers change \emph{correctly}, not merely that outputs differ. Public grounding results then connect this controlled behavior to practical localization performance.
